\subsection{The coefficient face of the same-field triples}

The construction also identifies how the three cycles enter the small
positive ratio sector. Its calibration varies $A,C$ and preserves the
three coefficients prescribed in \eqref{p4:path}. We compute the ratio
directly from those coefficients, independently of the upper estimates.

\begin{proposition}\label{prop:lower-coefficient-face}
The fields in Lemma~\ref{rew:lem-endpoint-attainment} can be chosen with
\[
 \alpha_0<0<\alpha_1,\qquad p\longrightarrow0^+,
 \qquad \gamma/s\longrightarrow+\infty.
\]
In their positive-gamma coordinates
$g=\gamma$, $k=s/g^{3/2}$, $\pi=p/\sqrt g$ and
$\mathfrak q=k\pi$, they satisfy
\[
 k\longrightarrow0,\qquad \pi\longrightarrow5/\sqrt2,
 \qquad \mathfrak q\longrightarrow0.
\]
Their two positive axis roots in the coordinate
$(v-\gamma)/(gk)$ have squared half-gap tending to $1/72$.
In particular they remain on the bounded positive-gamma real-root
side, separated from the fold. The same field and its three complete
hyperbolic cycles meet every prescribed positive upper cutoff on $p$.
\end{proposition}

\begin{proof}
For each fixed $c>1$, write
\[
 \epsilon=\left(\tfrac12e^{-4\pi/\sqrt{c-1}}\right)^{2/5},
 \qquad \beta=1-\tfrac52\epsilon.
\]
Restrict the fixed $c$-range so that $0<\epsilon<1/10$.
Equation~\eqref{p4:path}, with its fixed root ratio $3/2$, gives
\begin{equation}\label{eq:lower-sector-exact-coefficients}
 \gamma=\frac{\beta\rho}{3},\qquad
 \alpha_0=-\frac32\epsilon^2\rho^3,\qquad
 \alpha_1=\left(\frac52\epsilon-\frac32\epsilon^2\right)\rho^2.
\end{equation}
These are the actual coefficients after calibration: the equations
\eqref{p4:calibration} change only $A,C$. Moreover a field $H_\theta$
already in the slice has, by local uniqueness in Lemma~\ref{lem:slice},
slice parameter $\theta(H_\theta)=\theta$ and affine/time parameters
$(I,0,1)$. Consequently after the fixed inverse transport
$\mathscr C_0^{-1}$ these same coefficients define
\eqref{eq:sharp-sector-coordinates} for the original field.

Taking the square root and dividing in
\eqref{eq:lower-sector-exact-coefficients} gives
\begin{equation}\label{eq:lower-sector-ratios}
 s=\sqrt{3/2}\,\epsilon\rho^{3/2},\qquad
 p=\frac{5-3\epsilon}{\sqrt6}\sqrt\rho,\qquad
 \frac\gamma s=\frac{\beta}
 {3\sqrt{3/2}\,\epsilon\sqrt\rho}.
\end{equation}
For $c_n\downarrow1$, the endpoint construction first fixes $c_n$,
then chooses one $\rho_n\le1/n$ for the actual coefficient difference
and all three whole-cycle distances. Thus
$0<p_n\le5/\sqrt{6n}\to0$, while $\gamma_n/s_n\to+\infty$.
The actual coefficient $C_{c_n,\rho_n}$ tends to $c_n$ with the
chosen accuracy; it need not equal $c_n$. No estimate for the inverse
calibration uniform as $c\downarrow1$ is used.

The gamma coordinates computed from the same path are
\begin{equation}\label{eq:lower-sector-gamma-coordinates}
 k=\frac{9\epsilon}{\sqrt2\,\beta^{3/2}},\qquad
 \pi=\frac{5-3\epsilon}{\sqrt{2\beta}},\qquad
 \mathfrak q=\frac{9\epsilon(5-3\epsilon)}{2\beta^2}.
\end{equation}
The translated roots in \eqref{p4:path} are
$-\rho,\epsilon\rho,3\epsilon\rho/2$. Dividing by
$gk=3\epsilon\rho/\sqrt{2\beta}$ gives
\[
 -r=-\frac{\sqrt{2\beta}}{3\epsilon},\qquad
 a=\frac{\sqrt{2\beta}}3,\qquad b=\sqrt{\beta/2}.
\]
Hence
\begin{equation}\label{eq:lower-sector-root-gap}
 \frac{(b-a)^2}{4}=\frac{\beta}{72}\longrightarrow\frac1{72},
 \qquad kr=\frac3\beta\longrightarrow3.
\end{equation}
In particular $\pi\to5/\sqrt2>\sqrt{12}$, on the separated real-root
side, with an escaping negative root at finite further radius.
These formulas describe this constructed family; they impose no
gamma-chart limit on other three-cycle sequences.

For any neighborhood, collar and positive ratio cutoff, take a
sufficiently large $n$ in this same sequence. Coefficient convergence,
the simultaneous Hausdorff limits and $p_n\to0^+$ meet all three
requirements at once. The fixed inverse coordinate change and positive
time preserve the three hyperbolic cycles of that single field.
\end{proof}

The one-cycle estimate at $p=0$ and these strictly positive-ratio
triples concern different coefficient limits. The realization passes
through the unbounded $\gamma/s$ branch. A bounded-$\gamma/s$
restriction would remove this sharpness family.
